\documentclass{article}
\usepackage{graphicx} % Required for inserting images
\usepackage[letterpaper, portrait, margin=1in]{geometry}
\usepackage{float}
\usepackage{amsmath}
\usepackage{physics}
\usepackage{hyperref}
\usepackage{parskip} 

\title{Doctor, is this normal? \\ \large ULAB \\ Division of Sincerely Serious Science}
\author{Jonathan Zhang}
\date{\today}

\begin{document}

\maketitle

\section{Introduction}

According to my sources, around May 2021 an image began circulating around the internet about a broken sidewalk. The image, which can be seen in \autoref{fig:sidewalk}, claims that the sidewalk crack resembles the form of a normal distribution. This lab report will test that claim by fitting this curve with a standard normal distribution (also called a Gaussian distribution).

\begin{figure}[H]
    \centering
    \includegraphics[width=0.8\linewidth]{img.jpg}
    \caption{Sidewalk image from the internet.}
    \label{fig:sidewalk}
\end{figure}

\section{Methods}

\subsection{Preprocessing}

First, we used the Perspective Wrap tool in Photoshop to correct the tilted perspective of the image. Then, we worked with the \texttt{Python Imaging Library} and \texttt{NumPy} to convert our image to grayscale.

A threshold mask was applied to the image so that only dark pixels would be retained. Note that some minimal manual processing was required to remove the vertical sidewalk crack as well as outlier pixels.

Finally, the location of the dark pixels was recorded as coordinates in a Cartesian plane.

\subsection{Curve Fitting}

We worked with \texttt{SciPy} to fit a normal distribution:

\begin{equation}
f(x) = a \cdot e^{-\frac{1}{2} \left( \frac{x - \mu}{\sigma} \right)^2}
\label{eq:normal_distribution}
\end{equation}

where \(a\) is a scaling constant, \(\mu\) is the mean, and \(\sigma\) is the standard deviation.

\section{Results}

The results of the analysis are presented in \autoref{fig:fit}. This figure displays the data points extracted from the sidewalk crack and the fitted normal distribution, demonstrating the alignment between the observed data and the model.

We evaluated the fit using the reduced chi-squared statistic:

\begin{equation}
\chi^2_\nu = \frac{1}{\nu} \sum \frac{(O - E)^2}{E}
\label{eq:chi_squared}
\end{equation}

where \(\chi^2_\nu\) is the reduced chi-squared value, \(\nu\) is the number of degrees of freedom, \(O\) represents observed values, and \(E\) represents expected values.

Below is a table summarizing the curve fitting results:

\begin{table}[H]
\centering
\begin{tabular}{|c|c|}
\hline
\textbf{Parameter} & \textbf{Value} \\ \hline
Mean (\(\mu\))           & 0.50           \\ \hline
Standard Deviation (\(\sigma\)) & 0.10           \\ \hline
Reduced Chi-squared (\(\chi^2_\nu\)) & 0.22           \\ \hline
\end{tabular}
\caption{Summary of the curve fitting results for the sidewalk crack data.}
\label{tab:results}
\end{table}

\begin{figure}[H]
    \centering
    \includegraphics[width=13cm]{fit.png}
    \caption{Normal distribution fit to the sidewalk crack data. The curve represents the fitted Gaussian distribution overlaid on the observed data points.}
    \label{fig:fit}
\end{figure}

\bibliographystyle{ieeetr}  % Use IEEE style for references
\bibliography{ref.bib} 
\nocite{*}  % This will include all entries from the .bib file

\end{document}
